% ID: SEL-D0002-OA-E01
% ESTADO: validada. S003, 2026-11-07; tres pasadas del mismo autor, comprobación separada y revisión didáctica/visual.
\textbf{Datos y teoría.}\par\nopagebreak[4]
$A$ es $2\times3$; $A^t$ es $3\times2$, por lo que $AA^t$ es cuadrada de orden 2. Una matriz cuadrada es invertible si su determinante no es cero. En una ecuación matricial se conserva el orden de los factores.

\textbf{a) Simetría e inversa.}\par\nopagebreak[4]
Calculando productos escalares de las filas de $A$:
\[
 AA^t=\begin{pmatrix}5&-1\\-1&2\end{pmatrix}.
\]
La primera afirmación es \textbf{cierta}: coincide con su traspuesta. En general, $(AA^t)^t=(A^t)^tA^t=AA^t$.
\[
 AA^t+B=\begin{pmatrix}6&-4\\-3&2\end{pmatrix},\qquad
 \det(AA^t+B)=12-12=0.
\]
La segunda afirmación es \textbf{falsa}: esta matriz no tiene inversa.

\textbf{b) Resolución de $BX+A=C$.}\par\nopagebreak[4]
$X$ debe ser $2\times3$. Restamos $A$ y multiplicamos por $B^{-1}$ por la izquierda:
\[
 C-A=\begin{pmatrix}6&-12&18\\0&6&12\end{pmatrix},\quad
 \det B=-6,\quad
 B^{-1}=\begin{pmatrix}0&-1/2\\-1/3&-1/6\end{pmatrix}.
\]
\[
 \boxed{X=B^{-1}(C-A)=\begin{pmatrix}0&-3&-6\\-2&3&-8\end{pmatrix}.}
\]
Comprobación: $BX=\begin{pmatrix}6&-12&18\\0&6&12\end{pmatrix}=C-A$; al sumar $A$ se recupera exactamente $C$. La solución es única porque $B$ es invertible.
